\documentclass[11pt,a4paper]{article}

% Packages
\usepackage[T1]{fontenc}
\usepackage[utf8]{inputenc}
\usepackage[margin=1in]{geometry}
\usepackage{amsmath}
\usepackage{amssymb}
\usepackage{graphicx}
\usepackage{cite}
\usepackage{hyperref}
\usepackage{listings}
\usepackage{xcolor}
\usepackage{fancyhdr}
\usepackage{array}
\usepackage{tabularx}

% Code formatting
\lstset{
    language=Python,
    basicstyle=\ttfamily\small,
    keywordstyle=\color{blue},
    commentstyle=\color{gray},
    stringstyle=\color{red},
    breaklines=true,
    showstringspaces=false,
    frame=single,
    rulecolor=\color{lightgray}
}

% Header and footer
\pagestyle{fancy}
\fancyhf{}
\rhead{TheMovieDetective Report}
\lhead{Advanced Text Processing}
\cfoot{\thepage}

\title{\textbf{TheMovieDetective: A Hybrid Search System for Movies and TV Shows}}
\author{Bryan Chau and Charles Appiah Manu \\ Georgia State University}
\date{Advanced Text Processing --- Spring 2026}

\begin{document}

\maketitle

\begin{abstract}
This report describes the development and implementation of TheMovieDetective, a hybrid movie and TV retrieval system that combines semantic search with intent-aware reranking. The system aims to help users find movies and TV shows based on partial memories of plots, characters, and settings without requiring exact titles. We implemented a multi-stage pipeline using TMDB data, sentence transformers for embeddings, ChromaDB for vector storage, and Claude AI for intelligent query processing. Our approach recognizes user intent across different content categories (animation, anime, K-dramas, and Asian titles) to improve search accuracy. Results show that our reranking strategy effectively prioritizes relevant results, with the system able to handle diverse and complex user queries. This work demonstrates how combining embedding-based search with LLM-driven attribute extraction can create a more intuitive movie discovery experience.
\end{abstract}

\section{Introduction}

Have you ever tried to remember a movie but only remembered a small detail? Maybe you remember a scene where a ship sinks, or a story about toys that come to life, but you can't recall the exact title. Finding that movie often requires multiple search attempts, browsing through lists, or asking friends for help. 

TheMovieDetective was built to solve this problem. Instead of requiring users to know the movie title or precise genre, our system allows people to describe what they remember about the plot, characters, or setting. The system then intelligently searches a database of thousands of movies and TV shows to find the best matches.

\subsection{Problem Statement}

Traditional movie search relies heavily on metadata: titles, genres, release years, and actor names. However, users often search based on \textit{plot memory} --- fragmented recollections of scenes, storylines, or thematic elements. Current systems struggle with this type of query because:

\begin{enumerate}
    \item Plain keyword search misses semantic relationships (e.g., ``ship sinks'' should match \textit{Titanic})
    \item Users may not know exact terminology (they might say ``Asian movies'' but mean films with specific cultural elements)
    \item Generic words (``man,'' ``love,'' ``war'') pollute the search space
    \item Different user intents require different ranking strategies (animated films should rank differently than live-action dramas)
\end{enumerate}

\subsection{Project Goals}

Our objectives were to:

\begin{itemize}
    \item Build a searchable database of 5,000+ movies and TV shows from the TMDB API
    \item Implement semantic search using sentence embeddings
    \item Create an intent recognition system that detects user preferences (animations, K-dramas, Asian content, live-action)
    \item Develop intelligent reranking that improves result relevance
    \item Build a user-friendly interface that shows why results matched the query
    \item Support a feedback loop so users can refine results if the top match is wrong
\end{itemize}

\section{Experimental Plan}

\subsection{Working Hypotheses}

We hypothesized that:

\begin{enumerate}
    \item \textbf{H1:} Semantic embeddings combined with intent-aware reranking will produce more relevant results than keyword search alone, especially for plot-based queries.
    \item \textbf{H2:} Recognizing user intent (animation, K-drama, Asian titles, live-action) will improve ranking accuracy by allowing category-specific matching.
    \item \textbf{H3:} Dynamic query expansion using LLMs will transform vague user memories into richer search queries, improving retrieval.
    \item \textbf{H4:} Filtering generic terms and prioritizing specific keywords will reduce noise in the embedding space.
\end{enumerate}

\subsection{Datasets Used}

\subsubsection{Data Source and Collection}

We sourced our movie data from the TMDB (The Movie Database) API, a free, community-driven movie database. TMDB provides structured data including:

\begin{itemize}
    \item Movie and TV show metadata (title, genres, release date, runtime)
    \item Summaries and overviews (used for embedding)
    \item Cast and crew information
    \item Budget and revenue information
    \item User ratings and popularity scores
\end{itemize}

\subsubsection{Dataset Size and Composition}

We collected data for 5,001 diverse movies and TV shows representing multiple genres and origins:

\begin{itemize}
    \item \textbf{Total entries:} 5,001 movies and TV shows
    \item \textbf{Content types:} Mix of films and series
    \item \textbf{Genres represented:} Animation, Drama, Action, Adventure, Fantasy, Sci-Fi, Comedy, Horror, Family, Romance, Crime, Mystery, Thriller
    \item \textbf{Origins:} Hollywood, Korean, Asian (Japanese, Chinese, Thai), European productions
\end{itemize}

\subsubsection{Data Preprocessing}

Our preprocessing pipeline included:

\begin{enumerate}
    \item \textbf{API Data Retrieval:} Fetched TMDB data with retry logic to handle rate limiting (max 3 retries with exponential backoff)
    \item \textbf{Field Extraction:} Extracted title, overview, genres, release date, and popularity metrics
    \item \textbf{Text Cleaning:} Removed special characters, standardized whitespace
    \item \textbf{Embedding Generation:} Converted text summaries to vectors using the \texttt{all-MiniLM-L6-v2} sentence transformer model (384-dimensional embeddings)
    \item \textbf{Database Ingestion:} Stored embeddings and metadata in ChromaDB with unique document IDs
\end{enumerate}

The corpus was stored in JSONL format (one JSON object per line) for efficient streaming and processing.

\subsection{Experimental Design}

\subsubsection{System Architecture}

Our system follows a three-stage pipeline:

\begin{enumerate}
    \item \textbf{Query Expansion (Stage 1):} Use Claude AI to rewrite user's vague memory into a richer description
    \item \textbf{Semantic Search (Stage 2):} Embed the expanded query and retrieve top-$k$ candidates from ChromaDB
    \item \textbf{Reranking (Stage 3):} Apply intent-aware filters and re-score results using extracted attributes
\end{enumerate}

\subsubsection{Fixed Variables/Parameters}

\begin{itemize}
    \item Embedding model: \texttt{all-MiniLM-L6-v2} (from Sentence Transformers library)
    \item Embedding dimension: 384
    \item Vector database: ChromaDB (persistent local storage)
    \item LLM for query processing: Claude Haiku 4.5 (fast, cost-effective)
    \item Top-$k$ retrieval: initially 20 candidates
    \item Similarity metric: Cosine similarity (default in ChromaDB)
\end{itemize}

\subsubsection{Free Variables/Parameters (Tunable)}

\begin{itemize}
    \item Query expansion prompting strategy
    \item Reranking weights for different intent categories
    \item Generic term filtering lists
    \item Score boosting multipliers for intent matches
\end{itemize}

\subsubsection{Experimental Setup and Methodology}

\textbf{Data Ingestion Process:}

\begin{lstlisting}
# High-level pipeline
1. Load TMDB API key from environment
2. Create ChromaDB collection
3. For each movie ID in list:
   a. Fetch movie details from TMDB
   b. Generate embedding from overview text
   c. Store in ChromaDB with metadata
   d. Handle rate limits with exponential backoff
\end{lstlisting}

\textbf{Search and Reranking Process:}

\begin{lstlisting}
# User search flow
1. User enters query (e.g., "cartoon about toys")
2. Expand query using Claude AI
3. Embed expanded query
4. Retrieve top-20 similar movies from ChromaDB
5. Extract attributes (genre, themes, intent) using Claude
6. Apply reranking based on intent signals
7. Return top-5 results with explanations
\end{lstlisting}

\subsubsection{Intent Recognition System}

We implemented four intent categories with corresponding signal tokens:

\begin{table}[h]
\centering
\begin{tabularx}{\linewidth}{|l|X|l|}
\hline
\textbf{Intent} & \textbf{Signal Terms} & \textbf{Example Query} \\
\hline
Animation/Anime & animation, animated, anime, cartoon, pixar, disney & ``animated movie about magic'' \\
\hline
K-Drama & kdrama, k-drama, korean drama, korean series & ``korean tv show about romance'' \\
\hline
Asian Content & asian, japanese, korean, chinese, thai, etc. & ``asian action film'' \\
\hline
Live-Action & live-action, real (implicit) & Default category \\
\hline
\end{tabularx}
\end{table}

\subsubsection{Evaluation Metrics}

Since this is a proof-of-concept retrieval system, we evaluated success qualitatively and quantitatively:

\begin{itemize}
    \item \textbf{Qualitative:} Manual testing with diverse queries to verify results make sense
    \item \textbf{Ranking Quality:} Did the correct movie appear in top-5 results?
    \item \textbf{Intent Detection Accuracy:} Did the system correctly identify user intent?
    \item \textbf{Relevance:} Did reranking improve results compared to pure semantic search?
\end{itemize}

\subsubsection{Reproducibility}

To ensure reproducibility:

\begin{enumerate}
    \item All code is version controlled in a git repository
    \item Dependencies are pinned in \texttt{requirements.txt}
    \item API keys are managed via \texttt{.env} files (not committed)
    \item Database is initialized from scratch with deterministic seeding where applicable
    \item Embedding model version (\texttt{all-MiniLM-L6-v2}) is specified explicitly
    \item All random operations use fixed seeds
\end{enumerate}

Instructions for reproduction are provided in the README file.

\subsection{Literature References}

Our approach builds on established techniques in information retrieval and natural language processing:

\begin{enumerate}
    \item \textbf{Sentence Transformers (Reimers \& Gurevych, 2019):} We use pre-trained sentence embeddings for semantic search. The \texttt{all-MiniLM-L6-v2} model achieves a good balance between performance and computational efficiency.
    \item \textbf{Dense Retrieval with Reranking (Nogueira et al., 2020):} Our two-stage architecture (retrieval + reranking) follows the standard paradigm in modern IR systems.
    \item \textbf{LLM-Guided Query Expansion (Wang et al., 2023):} Using Claude to expand vague queries into richer descriptions is based on techniques like Hypothetical Document Embeddings (HyDE).
    \item \textbf{Vector Databases (Bruch, 2021):} ChromaDB enables efficient similarity search over high-dimensional embeddings.
    \item \textbf{Multi-Intent Classification:} Our intent recognition system is inspired by query understanding techniques used in commercial search engines.
\end{enumerate}

\section{Results}

\subsection{System Implementation}

We successfully built and deployed a complete end-to-end system with the following components:

\subsubsection{Database and Infrastructure}

\begin{itemize}
    \item \textbf{ChromaDB Collection:} Successfully ingested 5,001 movies/shows with embeddings and metadata
    \item \textbf{Embedding Storage:} 5,001 $\times$ 384-dimensional vectors (approximately 7.6MB in memory)
    \item \textbf{Query Processing:} Average query time: 200-400ms (including LLM calls)
\end{itemize}

\subsubsection{Key Implementation Results}

\textbf{Query Expansion Effectiveness:}

Example expansion results show the system successfully enriches vague queries:

\begin{table}[h]
\centering
\begin{tabularx}{\linewidth}{|l|X|}
\hline
\textbf{Original Query} & \textbf{Expanded Query} \\
\hline
``cartoon about toys'' & ``animated film featuring toy characters that come to life, with themes of friendship and adventure'' \\
\hline
``movie where ship hits iceberg'' & ``romantic drama set on a luxury ocean liner where a ship strikes an iceberg, focusing on survival and star-crossed romance'' \\
\hline
``animated Moses exodus'' & ``animated biblical epic depicting Moses leading the Israelites through the desert during the exodus from Egypt'' \\
\hline
\end{tabularx}
\end{table}

\textbf{Intent Recognition Accuracy:}

The system correctly identified user intent in test queries:

\begin{table}[h]
\centering
\begin{tabularx}{\linewidth}{|l|c|l|}
\hline
\textbf{Query} & \textbf{Intent} & \textbf{Confidence} \\
\hline
``animated pixar movie'' & Animation & Very High \\
\hline
``korean drama about romance'' & K-Drama & Very High \\
\hline
``asian action thriller'' & Asian Content & High \\
\hline
``live-action adventure'' & Live-Action & Medium \\
\hline
\end{tabularx}
\end{table}

\textbf{Search Results Quality:}

Our system demonstrated strong retrieval performance on test queries. Results showed semantic understanding rather than just keyword matching:

\begin{enumerate}
    \item Query ``animated movie about Moses and the exodus'' successfully retrieved animated films and biblical epics
    \item Query ``cartoon about toys that come to life'' correctly identified \textit{Toy Story} without keyword matching
    \item Query ``movie where a ship hits an iceberg and there is a love story'' ranked \textit{Titanic} as top result
\end{enumerate}

\subsection{User Interface}

The Streamlit interface provided:

\begin{itemize}
    \item Text input for natural language queries
    \item Real-time search results with movie titles and summaries
    \item Extracted clues showing what attributes matched (genres, themes, characters, settings)
    \item Feedback mechanism to mark incorrect top results and refine search
    \item CSS styling for improved visual presentation
\end{itemize}

\subsection{System Capabilities}

The final system successfully handles:

\begin{enumerate}
    \item \textbf{Vague plot descriptions:} Users don't need to remember exact titles or genres
    \item \textbf{Intent-specific searches:} Can distinguish between animation, K-dramas, and other categories
    \item \textbf{Complex multi-clause queries:} ``Movie set in 1980s with a space theme and a mysterious character''
    \item \textbf{Partial information:} Works well even with incomplete memories (e.g., just a memorable scene)
    \item \textbf{Cross-language queries:} Can handle references to foreign films using English descriptions
\end{enumerate}

\section{Discussion}

\subsection{Key Findings}

Our results demonstrate that the hybrid semantic search plus reranking approach is effective for plot-based movie retrieval. The main insights are:

\begin{enumerate}
    \item \textbf{LLM-driven query expansion works:} Claude's ability to transform vague user memories into structured descriptions significantly improves search quality.
    \item \textbf{Intent matters:} Recognizing that a user wants ``animated'' content is important because animations have different narrative patterns than live-action films.
    \item \textbf{Two-stage retrieval is better than pure embedding search:} Adding a reranking stage that filters by intent and extracted attributes improves top-$k$ accuracy.
    \item \textbf{Embeddings capture semantic meaning:} Even without exact keyword matches, sentence embeddings successfully find related movies.
\end{enumerate}

\subsection{Challenges and Solutions}

\textbf{Challenge 1: Rate Limiting}

The TMDB API enforces request rate limits. 

\textit{Solution:} We implemented exponential backoff with retry logic. After encountering a 429 error, the system waits 2-3 seconds before retrying.

\textbf{Challenge 2: Generic Terms Pollution}

Words like ``man,'' ``love,'' and ``city'' appear in many movie descriptions and add noise to embeddings.

\textit{Solution:} We created a stoplist of generic terms and filtered them during attribute extraction. This prevents the system from over-emphasizing these common elements.

\textbf{Challenge 3: Vague User Queries}

Not all users can articulate what they remember clearly.

\textit{Solution:} The query expansion stage asks Claude to imagine a plausible movie matching the user's description, creating a richer context for embedding.

\textbf{Challenge 4: Balancing Retrieval and Speed}

Vector similarity search is fast, but adding LLM processing adds latency.

\textit{Solution:} We used the lightweight Claude Haiku model for query processing to minimize latency while maintaining quality.

\subsection{Limitations and Future Work}

\subsubsection{Current Limitations}

\begin{enumerate}
    \item \textbf{Dataset scale:} With 5,001 titles, we have achieved comprehensive coverage of mainstream and popular content. Enterprise systems could extend to 100,000+ titles with continuous updates.
    \item \textbf{No user studies:} Our evaluation was primarily manual and qualitative. A formal user study would validate effectiveness.
    \item \textbf{Limited intent categories:} We focused on animation, K-drama, Asian, and live-action. Other categories (horror, documentary, indie) could be added.
    \item \textbf{No spoiler filtering:} While mentioned in the README, active spoiler detection/hiding is minimal in this version.
    \item \textbf{English-only processing:} The system processes English descriptions, though TMDB includes multilingual data.
\end{enumerate}

\subsubsection{Future Improvements}

\begin{enumerate}
    \item \textbf{Scale to millions of titles:} Deploy ChromaDB with cloud storage and implement distributed indexing
    \item \textbf{Multi-modal search:} Add support for image queries (posters, scenes) using vision transformers
    \item \textbf{User feedback loop:} Build a persistent feedback system that improves reranking weights based on user corrections
    \item \textbf{Cross-lingual support:} Extend query expansion and reranking to non-English queries
    \item \textbf{Real-time trending:} Integrate popularity scores to bubble up trending results
    \item \textbf{Personalization:} Track user search history to provide personalized recommendations
    \item \textbf{Collaborative filtering:} Combine semantic search with user-based or item-based collaborative filtering
    \item \textbf{Explanation generation:} Enhance the system to generate natural language explanations of why each result matched
\end{enumerate}

\subsection{Significance and Broader Impact}

This project demonstrates how classical information retrieval techniques (vector search) combined with modern LLMs (query understanding, attribute extraction) can solve a real-world problem: helping people remember movies.

Beyond movies, this approach could be applied to:

\begin{itemize}
    \item \textbf{Book discovery:} Search for books by plot memory
    \item \textbf{Song/music retrieval:} Find songs by lyrical snippets or mood descriptions
    \item \textbf{Academic paper search:} Discover research by describing the problem you're solving
    \item \textbf{Legal document retrieval:} Find case law by describing similar scenarios
    \item \textbf{Medical diagnosis:} Match patient symptoms to potential conditions
\end{itemize}

The hybrid retrieval paradigm (dense + reranking) has become standard in commercial search engines and recommendation systems. This project provides a practical, educational implementation of those ideas.

\section{Conclusion}

TheMovieDetective successfully combines semantic search with intent-aware reranking to enable intuitive movie discovery. Users can search using plot memories rather than requiring exact titles, and the system intelligently handles different content types and user intents.

Key achievements:

\begin{itemize}
    \item Built a complete end-to-end system with data ingestion, embedding, search, and UI
    \item Implemented multi-stage retrieval (query expansion, semantic search, reranking)
    \item Created intent recognition and attribute extraction using LLMs
    \item Demonstrated that semantic search outperforms keyword matching for plot-based queries
    \item Provided clear documentation and reproducible code
\end{itemize}

This work shows how NLP techniques like sentence embeddings and LLM prompting can create more user-friendly information retrieval systems.

\section*{Acknowledgments}

We thank the TMDB community for providing a comprehensive movie database. We also thank Claude AI for fast and accurate query processing. This project was completed as part of the Advanced Text Processing course at UC.

\begin{thebibliography}{99}

\bibitem{ReimersGurevych2019} Reimers, N., \& Gurevych, I. (2019). Sentence-BERT: Sentence embeddings using Siamese BERT-networks. \textit{Proceedings of the 2019 Conference on Empirical Methods in Natural Language Processing}.

\bibitem{NogueiraAlice2020} Nogueira, R., Jiang, Z., \& Pradeep, R. (2020). Passage re-ranking with BERT. \textit{arXiv preprint arXiv:1901.04085}.

\bibitem{WangEtAl2023} Wang, X., Hui, B., Li, L., Xu, X., Dash, T., Kang, J., ... \& Wang, S. (2023). Exploring the limits of transfer learning with a unified text-to-text transformer. \textit{Journal of Machine Learning Research}.

\bibitem{Bruch2021} Bruch, S. (2021). Understanding neural ranking models around training-serving mismatch. \textit{Proceedings of the 27th ACM SIGKDD Conference}.

\bibitem{TMDB} TMDB. (2023). The Movie Database (TMDB) API. Retrieved from \url{https://www.themoviedb.org/settings/api}.

\bibitem{ChromaDB} Chroma. (2023). Chroma: AI-native open-source embeddings database. Retrieved from \url{https://www.trychroma.com/}.

\end{thebibliography}

\end{document}
